\ifwithsols
\else
\clearpage
\fi

\subsection*{سؤال 2 امتحان 899371 سنة 2025}
\addcontentsline{toc}{subsection}{سؤال 2 امتحان 899371 سنة 2025}

\noindent
\makebox[\textwidth][c]{\includegraphics[width=0.9\paperwidth,keepaspectratio]{ ../../../bagrut_questions/basics/arrays_2025_899371_2.jpg }%
}%
\begin{boxHint}
    استعمل مصفوفة مساعدة، انسخ عليها أولا الأعداد الموجبة، ثم الأعداد السالبة، ثم انسخها إلى المصفوفة الأصلية.
\end{boxHint}

\ifwithsols
\begin{boxSolution}[سؤال 2 امتحان 899371 سنة 2025]
هناك عدة طرق لحل هذا السؤال:
\begin{enumerate}
    \item الطريقة الأولى: استخدام مصفوفة واحدة مساعدة
    \item الطريقة الثانية: استخدام مصفوفتين مساعدتين: واحدة لتخزين الأعداد الموجبة، والثانية للأعداد السالبة، ثم ننقلها بالترتيب إلى المصفوفة الأصلية.
    \item الطريقة الثالثة: استخدام مصفوفة واحدة لتخزين الأعداد السالبة فقط، ثم ننقل الأعداد الموجبة إلى بداية المصفوفة الأصلية، ثم ننسخ الأعداد السالبة من المصفوفة المساعدة إلى نهاية المصفوفة الأصلية.
    \item الطريقة الرابعة: بدون استخدام مصفوفة مساعدة: نبدأ من نهاية المصفوفة، في كل مرة نجد عددًا سالبًا ننسخه إلى نهاية المصفوفة (قبل الأعداد السالبة التي وجدناها حتى الآن) وننفذ إزاحة للأعداد الموجبة إلى اليسار.
\end{enumerate}
سنعرض حل الطريقة الأولى والرابعة فقط - الثانية والثالثة يمكنك كتابتها لوحدك (تشبه الطريقة الأولى).
\end{boxSolution}
\clearpage
\begin{boxSolution}[الطريقة الأولى: استخدام مصفوفة مساعدة]
{\small
\begin{itemize}
    \item في الحلقة الأولى، ننسخ جميع الأعداد الموجبة إلى مصفوفة مساعدة.
    \item في الحلقة الثانية، ننسخ جميع الأعداد السالبة إلى نفس المصفوفة المساعدة، بحيث تكون بعد الأعداد الموجبة.
    \item في الحلقة الثالثة، ننسخ العناصر من المصفوفة المساعدة إلى المصفوفة الأصلية.
\end{itemize}
}
\begin{boxCode}
\begin{english}
\begin{minted}{csharp}
public static void Arrange(int[] arr)
{
    int[] temp = new int[arr.Length];
    int index = 0;
    for (int i = 0; i < arr.Length; i++)
    {
        if (arr[i] >= 0)
        {
            temp[index] = arr[i];
            index++;
        }
    }
    for (int i = 0; i < arr.Length; i++)
    {
        if (arr[i] < 0)
        {
            temp[index] = arr[i];
            index++;
        }
    }
    for (int i = 0; i < arr.Length; i++)
    {
        arr[i] = temp[i];
    }
}
\end{minted}
\end{english}
\end{boxCode}
\end{boxSolution}

\clearpage

\begin{boxSolution}[الطريقة الرابعة: بدون استخدام مصفوفة مساعدة]
\begin{boxCode}
\begin{english}
\begin{minted}{csharp}
public static void Arrange(int[] arr)
{
    int negIdx = arr.Length - 1;
    //int nextNegIdx = arr.Length - 1;
    for (int i = arr.Length - 1; i >= 0; i--)
    {
        if (arr[i] < 0)
        {
            if (i == negIdx)
                negIdx--;
            else if (i < negIdx)
            {
                int tmp = arr[i];
                // shift all positive numbers between i and negIdx to the left
                for (int j = i; j < negIdx; j++)
                    arr[j] = arr[j + 1];
                arr[negIdx] = tmp;
                negIdx--;
            }
        }
    }
}
\end{minted}
\end{english}
\end{boxCode}

\end{boxSolution}
\clearpage
\fi
